\documentclass[11pt]{article}
\usepackage[margin=1in]{geometry}
\usepackage{enumitem}
% =============================================================================
% Preambles/header.tex — shared LaTeX/Beamer preamble for this project
%
% Use from a Beamer lecture by prepending:
%     \input{header}
% and compiling with TEXINPUTS=../Preambles:$TEXINPUTS (the /compile-latex
% skill does this automatically).
%
% PALETTE CONTRACT. Color names defined here MUST match the names in
% Quarto/theme-template.scss so Beamer and Quarto renderings use the same
% palette. The scripts/check-palette-sync.sh script enforces this.
%
% To customize: edit the HEX values below AND the matching $variable in
% Quarto/theme-template.scss, then run scripts/check-palette-sync.sh.
% =============================================================================

% -----------------------------------------------------------------------------
% Engine detection + encoding
%
% Our /compile-latex skill uses XeLaTeX, where inputenc is unnecessary and
% can produce encoding warnings. Only load inputenc under pdfTeX.
% -----------------------------------------------------------------------------
\usepackage{iftex}
\ifPDFTeX
  \usepackage[utf8]{inputenc}
\fi

\usepackage{amsmath, amssymb, amsthm}
\usepackage{booktabs}
\usepackage{graphicx}

% -----------------------------------------------------------------------------
% PALETTE — mirrors Quarto/theme-template.scss variable names.
% Load xcolor and define colors BEFORE anything that references them
% (notably hyperref, hypersetup, and the Beamer color assignments).
% -----------------------------------------------------------------------------
\usepackage{xcolor}

\definecolor{primary-blue}{HTML}{012169}      % headings, accents
\definecolor{primary-gold}{HTML}{B9975B}      % emphasis, borders
\definecolor{highlight-yellow}{HTML}{F2A900}  % markers, alerts
\definecolor{light-bg}{HTML}{E8EDF5}          % title / section backgrounds
\definecolor{jet}{HTML}{1A1A1A}               % body text

% Semantic colors (used in diagrams and annotations)
\definecolor{positive}{HTML}{15803D}          % good / observed / identified
\definecolor{negative}{HTML}{B91C1C}          % bad / problematic / bias
\definecolor{neutral}{HTML}{525252}           % reference / context

% Highlight accents (match .hi-* classes in the SCSS)
\definecolor{hi-slate}{HTML}{314F4F}
\definecolor{hi-green}{HTML}{2E8B57}
\definecolor{hi-red}{HTML}{C41E3A}

% Semantic aliases so skill templates and snippets can write
% \textcolor{accent}{...} without hard-coding "primary-blue".
\colorlet{accent}{primary-blue}
\colorlet{accent2}{primary-gold}

% -----------------------------------------------------------------------------
% Hyperref — loaded AFTER xcolor + palette so link colors resolve.
% -----------------------------------------------------------------------------
\usepackage{hyperref}
\hypersetup{colorlinks=true, linkcolor=primary-blue, urlcolor=primary-blue,
            citecolor=primary-blue, pdfborder={0 0 0}}

% -----------------------------------------------------------------------------
% Course code — set once per deck (after \input{header}, before \begin{document})
% via \coursecode{CS401}. Shown in the footer on every slide so a deck's course
% is identifiable without opening the title page. Empty by default.
% -----------------------------------------------------------------------------
\makeatletter
\newcommand{\coursecode}[1]{\gdef\@coursecodetext{#1}}
\gdef\@coursecodetext{}
\makeatother

% -----------------------------------------------------------------------------
% Beamer theme assignments (only apply under Beamer).
%
% We detect Beamer via \@ifclassloaded{beamer}, which is the documented,
% non-internal way. This must be wrapped in \makeatletter/\makeatother
% because \@ifclassloaded contains an @-catcode character.
% -----------------------------------------------------------------------------
\makeatletter
\@ifclassloaded{beamer}{%
  \usefonttheme{professionalfonts}
  \setbeamercolor{structure}{fg=primary-blue}
  \setbeamercolor{frametitle}{fg=primary-blue}
  \setbeamercolor{title}{fg=primary-blue}
  \setbeamercolor{subtitle}{fg=primary-gold}
  \setbeamercolor{author}{fg=primary-blue}
  \setbeamercolor{institute}{fg=neutral}
  \setbeamercolor{date}{fg=primary-gold}
  \setbeamercolor{itemize item}{fg=primary-blue}
  \setbeamercolor{itemize subitem}{fg=primary-gold}
  \setbeamercolor{alerted text}{fg=primary-gold}
  \setbeamercolor{block title}{fg=white, bg=primary-blue}
  \setbeamercolor{block body}{bg=light-bg}
  % Semantic block variants: block=definition/concept (blue), exampleblock=
  % worked example (green/positive), alertblock=key takeaway or caution (gold).
  % Keep to <=2 colored boxes per slide across ALL three variants (INV-7).
  \setbeamercolor{block title example}{fg=white, bg=positive}
  \setbeamercolor{block body example}{bg=light-bg}
  \setbeamercolor{block title alerted}{fg=white, bg=primary-gold}
  \setbeamercolor{block body alerted}{bg=light-bg}

  % Minimal footer: course code (left) + page X / N (right), no navigation clutter
  \setbeamertemplate{navigation symbols}{}
  \setbeamertemplate{footline}{%
    \color{neutral}\footnotesize\hspace{0.7em}\@coursecodetext\hfill
    \insertframenumber/\inserttotalframenumber\hspace{0.7em}\vspace{0.3em}}
}{}
\makeatother

% -----------------------------------------------------------------------------
% TikZ setup — libraries the snippet gallery uses
% -----------------------------------------------------------------------------
\usepackage{tikz}
\usetikzlibrary{arrows.meta, positioning, calc, decorations.pathreplacing,
                fit, shapes.geometric, backgrounds}

% Shared TikZ styles — snippets and hand-written diagrams can pick these up
% via `\begin{tikzpicture}[every node/.style={...}]` or by naming specific
% styles (e.g., `\node[dag-node] (X) at (0,0) {X};`).
\tikzset{
  % Node styles for causal DAGs and similar
  dag-node/.style = {
    draw=primary-blue, fill=light-bg, rounded corners=2pt,
    minimum width=1.2cm, minimum height=0.9cm, align=center, font=\small
  },
  decision-node/.style = {
    draw=primary-gold, fill=white, diamond, aspect=1.8,
    minimum width=1.8cm, minimum height=1.2cm, align=center, font=\small,
    inner sep=1pt
  },
  % Edge styles — observed vs counterfactual
  observed-edge/.style = {-{Latex[length=2.5mm]}, thick, draw=jet},
  counterfactual-edge/.style = {-{Latex[length=2.5mm]}, thick, draw=neutral, dashed},
  confound-edge/.style = {-{Latex[length=2.5mm]}, thick, draw=negative, dashed},
  % Data-point styles for plots
  observed-dot/.style = {fill=primary-blue, draw=primary-blue, circle, inner sep=1.2pt},
  counterfactual-dot/.style = {fill=white, draw=neutral, circle, inner sep=1.2pt}
}

% -----------------------------------------------------------------------------
% Convenience macros
% -----------------------------------------------------------------------------
\newcommand{\muted}[1]{\textcolor{neutral}{#1}}
\newcommand{\key}[1]{\textcolor{primary-gold}{\textbf{#1}}}
\newcommand{\good}[1]{\textcolor{positive}{#1}}
\newcommand{\bad}[1]{\textcolor{negative}{#1}}

% Standout frame macro: full-bleed dark background for section transitions.
% Beamer-only (uses \begin{frame}/\setbeamercolor/\usebeamerfont) -- do not
% call from an article-class file (e.g. Notes/<CODE>/*-notes.tex). \muted,
% \key, \good, \bad, and \coursecode above are plain xcolor/gdef and are
% safe to use from either class; verified when Notes/article-class support
% was added.
% Expand: \transitionslide{Your transition title}
\newcommand{\transitionslide}[1]{%
  \begingroup
  \setbeamercolor{background canvas}{bg=primary-blue}
  \begin{frame}[plain]
    \centering
    \vfill
    {\usebeamerfont{title}\color{white}\Large #1\par}
    \vfill
  \end{frame}
  \endgroup
}

% Section divider frame (Beamer-only), an alternative to \transitionslide with a
% darker two-line style: near-black (jet) background, a small white label line
% above a larger gold title, and a thin gold rule along the bottom edge.
% Expand: \sectiondivider{Section 2}{Pointers}
\newcommand{\sectiondivider}[2]{%
  \begingroup
  \setbeamercolor{background canvas}{bg=jet}
  \begin{frame}[plain]
    \vspace*{3.0cm}
    {\color{white}\ttfamily\large #1\par}
    \vspace{0.5cm}
    {\color{primary-gold}\LARGE #2\par}
    \begin{tikzpicture}[remember picture, overlay]
      \draw[primary-gold, line width=2.5pt]
        ($(current page.south west)+(0,0.1)$) -- ($(current page.south east)+(0,0.1)$);
    \end{tikzpicture}
  \end{frame}
  \endgroup
}

% -----------------------------------------------------------------------------
% Channel watermark — "Concepts That Click" (YouTube).
%
% Article-class files (CompetitiveExam Q/A sets, Notes, Assignments) get a
% light diagonal draft watermark on every page plus a centered footer naming
% the channel. Beamer decks get a subtle TikZ background watermark instead
% (the footline already carries the course code). Centralized here so every
% MCQ PDF is branded without per-file edits.
% -----------------------------------------------------------------------------
\makeatletter
\@ifclassloaded{article}{%
  \usepackage{draftwatermark}
  \SetWatermarkText{Concepts That Click}
  \SetWatermarkScale{1}
  \SetWatermarkFontSize{2cm}
  \SetWatermarkLightness{0.86}
  \SetWatermarkAngle{45}
  \usepackage{fancyhdr}
  \pagestyle{fancy}
  \fancyhf{}
  \fancyfoot[C]{\footnotesize\textcolor{neutral}{Concepts That Click~|~YouTube: \href{https://www.youtube.com/@concepts.thatclick}{@concepts.thatclick}}}
  \renewcommand{\headrulewidth}{0pt}
  \renewcommand{\footrulewidth}{0pt}
  \fancypagestyle{plain}{%
    \fancyhf{}%
    \fancyfoot[C]{\footnotesize\textcolor{neutral}{Concepts That Click~|~YouTube: \href{https://www.youtube.com/@concepts.thatclick}{@concepts.thatclick}}}%
    \renewcommand{\headrulewidth}{0pt}%
    \renewcommand{\footrulewidth}{0pt}%
  }
}{}
\@ifclassloaded{beamer}{%
  \setbeamertemplate{background}{%
    \begin{tikzpicture}[remember picture, overlay]
      \node[rotate=45, opacity=0.055, align=center]
        at (current page.center)
        {\Huge\textcolor{neutral}{Concepts That Click}};
    \end{tikzpicture}%
  }
}{}
\makeatother 
\coursecode{JKSSB}
\renewcommand{\thesection}{45.\arabic{section}}
\newtheorem{example}{Example}[section]
\newenvironment{definitionbox}[1]{%
  \par\noindent\textbf{Definition (#1).}\ \itshape
}{\par\vspace{0.3em}}
\title{Open Source --- Detailed Lecture Notes}
\author{JKSSB Computer Awareness}
\date{\today}

\begin{document}
\maketitle

\section{Source code and what it means}
Every program begins as \textbf{source code}: the set of instructions a
programmer writes in a programming language such as C, Java, or Python. Source
code is \emph{human-readable} --- a person can open the file, read it, and
change it. Before the computer can run the program, a \textbf{compiler} or an
\textbf{interpreter} converts the source code into a machine-readable form.
That machine-readable form is called \textbf{object code}, or simply the
executable, and it is what the processor actually executes. Ordinary users of
a finished program normally receive only this compiled form and never see the
source code.

\begin{definitionbox}{Source code}
The human-readable set of instructions, written in a programming language,
that defines what a program does; it is converted into machine-readable
object code by a compiler or an interpreter.
\end{definitionbox}

The important consequence is that whoever holds the source code controls how
the software can be studied, changed, and shared. The whole topic of open
source versus proprietary software grows out of this single fact.

\section{Open source versus proprietary software}
Software falls into two broad camps according to what is done with its source
code.

\begin{center}
\begin{tabular}{p{0.22\textwidth}p{0.34\textwidth}p{0.34\textwidth}}
\toprule
\textbf{Point} & \textbf{Open source} & \textbf{Proprietary (closed source)} \\
\midrule
Source code & Publicly available & Kept secret \\
View or modify & Allowed, subject to the licence & Not allowed without permission \\
Redistribute & Allowed, subject to the licence & Not allowed \\
Typical user gets & Source code plus executable & Only the compiled executable \\
\bottomrule
\end{tabular}
\end{center}

\textbf{Open-source software} makes its source code public. Anyone is allowed
to inspect it, modify it, and redistribute it, subject to the terms of its
licence. \textbf{Proprietary software} keeps its source code confidential; the
user receives only the compiled executable and may not view, change, or share
the source without the owner's permission. Examples of proprietary software
include Microsoft Windows, Apple macOS, Microsoft Office, Adobe Photoshop, and
Oracle Database. The deciding factor between the two is the source code, not
the price: some open-source software is sold, and some proprietary software is
given away free of cost.

\begin{definitionbox}{Open source}
A software development model in which the source code is made publicly
available so that anyone may view, modify, and redistribute it under the
terms of its licence.
\end{definitionbox}

\section{The open-source idea: view, modify, share}
The meaning that examiners most often test is the three-part permission of an
open-source licence. First, a user may \textbf{view} the source code and see
exactly how the software works. Second, a user may \textbf{modify} it --- fix a
bug, improve it, or add a new feature. Third, a user may \textbf{share} it,
redistributing copies and modified versions to others. These permissions are
granted by the software's licence, so the exact conditions differ from one
project to another; but the underlying spirit is the same three freedoms in
the code.

\begin{example}
An item states, ``An open-source licence allows a user to view, change, and
share the source code.'' This is the standard, correct description and is the
answer expected in the exam. An option that says only ``software available
free of cost'' describes \emph{freeware}, not open source.
\end{example}

\section{Freeware, free software, public domain, commercial, shareware}
Several labels are often mixed up with open source. Each has a distinct
meaning.

\begin{center}
\begin{tabular}{p{0.20\textwidth}p{0.34\textwidth}p{0.36\textwidth}}
\toprule
\textbf{Term} & \textbf{What is free} & \textbf{Source code} \\
\midrule
Open source & Use, study, change, share & Available \\
Freeware & Cost only & Usually not provided \\
Free software & The user's freedom & Available \\
Shareware & A limited trial period & Not provided \\
Public domain & No copyright at all & No restrictions \\
Commercial & Nothing (it is sold) & Usually not provided \\
\bottomrule
\end{tabular}
\end{center}

\textbf{Freeware} is software given away free of cost, but its source code is
usually not provided. \textbf{Free software} is about freedom, not price: the
user is free to study, change, and share it. \textbf{Public domain} software
has no copyright owner, so anyone may use, modify, and redistribute it without
any restriction. \textbf{Shareware} is distributed as a free trial, working
for a limited period or with limited features; to keep using it, the user must
pay. \textbf{Commercial software} is developed and sold for profit and is
often, though not always, proprietary.

\begin{definitionbox}{Freeware}
Software distributed free of cost, whose source code is usually not made
available; it is free in price, not in freedom.
\end{definitionbox}

\begin{definitionbox}{Shareware}
Software distributed as a trial version, free for a limited time or with
limited features, after which payment is required to continue using it.
\end{definitionbox}

\section{Free software and its four freedoms}
The term \textbf{free software} comes from the \textbf{Free Software
Foundation (FSF)}, founded by \textbf{Richard Stallman}. In this context,
``free'' refers to liberty, not price --- the well-known phrase is ``free as in
freedom, not free as in free beer''. The FSF defines four freedoms: the freedom
to run the program for any purpose; the freedom to study how it works and
change it; the freedom to redistribute copies to help others; and the freedom
to distribute modified versions to the public. Free software and open source
are closely related and frequently overlap, but they are not identical labels.
Free software stresses the ethical freedom of the user; open source stresses
the practical benefits of a publicly available source code.

\begin{definitionbox}{Free software}
Software that grants the user the freedom to run, study, change, and
redistribute it and its modified versions; ``free'' here means freedom, not
zero cost.
\end{definitionbox}

\section{The Open Source Initiative}
Alongside the Free Software Foundation stands the \textbf{Open Source
Initiative (OSI)}, founded in 1998. The OSI maintains the \textbf{Open Source
Definition}, the set of criteria a licence must satisfy to be called an
open-source licence, and it reviews and approves licences against those
criteria. The term \textbf{open source} itself was popularised in 1998 as a
more business-friendly label than ``free software''. The two movements
overlap heavily: most licences approved by the OSI are also free-software
licences. An umbrella term often seen in questions is \textbf{FOSS}, meaning
Free and Open-Source Software, or \textbf{FLOSS} (Free/Libre and Open-Source
Software), which covers both traditions together.

The Open Source Definition lays down the conditions a licence must meet. The
software must be freely redistributable; the source code must be included or
easily obtainable; derived works and modifications must be allowed; the
licence must not discriminate against any person, group, or field of use; and
the rights must apply to everyone who receives the software. A licence that
satisfies these conditions is called an \textbf{OSI-approved} licence. This is
why ``open source'' is a matter of the licence, not simply of whether the
software is free of cost.

\begin{example}
Two programs are both free to download. Program A ships with its source code
under an OSI-approved licence that permits modification and redistribution;
Program B ships only compiled files and forbids copying. A is open source; B
is \emph{freeware}. The download price is identical, but the licences differ,
and the licence is what decides the category.
\end{example}

\section{Benefits and limitations of open source}
Open source is chosen for several practical reasons. The licence cost is
usually zero, so organisations avoid per-seat fees. The source code can be
inspected, which lets many people review it for bugs and security weaknesses.
Users are not locked in to a single vendor, and the software can be
customised to local needs. A large community often provides documentation and
support.

Open source also has limitations that a fair answer should acknowledge. A
popular project may have no formal commercial support contract, so support
depends on the community or on paid third-party vendors. Compatibility with
proprietary formats and integration with existing systems can be difficult.
The total cost of ownership is not always zero, because training, migration,
and maintenance still cost money. These points are occasionally tested as
``which of the following is a limitation of open source''.

\section{Full forms and key names to remember}
Direct-recall items often ask for an expansion or for the name behind a term.
The following are the ones most worth memorising.

\begin{center}
\begin{tabular}{p{0.22\textwidth}p{0.68\textwidth}}
\toprule
\textbf{Short form} & \textbf{Full form / meaning} \\
\midrule
GPL & GNU General Public License \\
FSF & Free Software Foundation \\
OSI & Open Source Initiative \\
FOSS & Free and Open-Source Software \\
FLOSS & Free/Libre and Open-Source Software \\
AOSP & Android Open Source Project \\
MIT & MIT License (a permissive software licence) \\
BSD & Berkeley Software Distribution (its licence family) \\
\bottomrule
\end{tabular}
\end{center}

Remember too the names behind the ideas: the Free Software Foundation was
founded by \textbf{Richard Stallman}, and \textbf{Git} was created by Linus
Torvalds to manage the Linux kernel's source code. \textbf{GitHub},
\textbf{GitLab}, \textbf{Bitbucket}, and \textbf{SourceForge} are the
best-known hosting sites for open-source repositories.

\section{Repositories and community development}
Open-source projects live in a \textbf{repository} (often shortened to
``repo''). A repository is a storage location that holds a project's source
code together with every change ever made to it. \textbf{Version control}
records this history, so an earlier version can be restored if needed.
\textbf{Git} is the most widely used distributed version-control system, and
popular hosting sites for repositories include \textbf{GitHub}, \textbf{GitLab},
\textbf{Bitbucket}, and \textbf{SourceForge}.

Open-source software is usually built by a \textbf{community} of contributors
spread around the world. A contributor may \textbf{fork} a project, that is,
make a personal copy for experimenting or improving. A \textbf{pull request}
is then a proposal to merge those changes back into the main project, where
\textbf{maintainers} review them. An \textbf{issue tracker} records bugs and
feature requests. Because many people can inspect and improve the code,
community review is treated as a key strength of open-source development.

\section{Basic licensing concepts}
By default, \textbf{copyright} gives the creator of a work exclusive rights
over it. A \textbf{licence} is the legal permission that states what others
may do with the software. Open-source and free-software licences fall into two
broad families.

\textbf{Copyleft} licences, such as the GNU General Public License (GPL),
require that modified versions of the software be released under the
\emph{same} licence. This keeps the software and its derivatives open.
\textbf{Permissive} licences, such as the MIT License, the Apache License 2.0,
and the BSD License, allow the code to be reused with very few conditions,
including inside proprietary projects, provided the licence notice is kept.
The licence, not the price, decides whether software is truly open source.

\begin{definitionbox}{Copyleft}
A licence condition that requires any modified or extended version of the
software to be distributed under the same licence, so the work cannot be
closed.
\end{definitionbox}

\section{Copyright, licences, and piracy}
It helps to separate three ideas. \textbf{Copyright} is the creator's legal
ownership of the work. A \textbf{licence} is the permission the owner grants
for specific uses. \textbf{Piracy} is the use, copying, or distribution of
software in a way the licence forbids --- for example, installing a
proprietary program without a valid licence, or copying it to machines the
licence does not cover. Open-source software is not ``unlicensed'': it is
licensed, and its licence grants broad permissions. Public-domain software,
by contrast, has no owner and therefore needs no licence at all. The exam
sometimes tests this by asking whether open-source software is ``free of
copyright''; the answer is no, because the copyright remains with the authors
and the licence grants the freedoms.

\begin{example}
A student installs the same licensed copy of a proprietary office suite on
five computers when the licence permits only one. This is \textbf{piracy}.
Installing an open-source office suite such as LibreOffice on all five
computers, in line with its licence, is not piracy. The difference is not the
price paid but whether the use matches what the licence allows.
\end{example}

\section{Common licences at a glance}
\begin{center}
\begin{tabular}{p{0.20\textwidth}p{0.30\textwidth}p{0.40\textwidth}}
\toprule
\textbf{Licence} & \textbf{Type} & \textbf{Main condition} \\
\midrule
GNU GPL & Copyleft & Derivative works must use the same licence \\
MIT License & Permissive & Reuse allowed; keep the licence notice \\
Apache License 2.0 & Permissive & Few limits, plus an explicit patent grant \\
BSD License & Permissive & Very few conditions \\
Creative Commons & Content licence & For creative works, not usually software \\
\bottomrule
\end{tabular}
\end{center}

Note that \textbf{Creative Commons} licences are designed for creative works
such as text, images, and music; they are not typically used to license
software. For software, the GPL, MIT, Apache, and BSD licences are the ones to
recognise.

\section{Well-known open-source and proprietary software}
Familiar examples help fix the concepts.

\begin{center}
\begin{tabular}{p{0.46\textwidth}p{0.46\textwidth}}
\toprule
\textbf{Open source} & \textbf{Proprietary} \\
\midrule
Linux (operating system) & Microsoft Windows \\
Android (mobile operating system) & Apple macOS, iOS \\
Mozilla Firefox (browser) & Microsoft Office \\
Chromium (Chrome's open-source base) & Adobe Photoshop \\
LibreOffice (office suite) & Oracle Database \\
VLC media player, GIMP, Blender & \\
Apache HTTP Server, MySQL, Python, Git, WordPress & \\
\bottomrule
\end{tabular}
\end{center}

\textbf{Android} is the open-source \textbf{mobile operating system}, built on
the open-source Linux kernel. This is a directly tested fact: when a question
asks for the open-source mobile OS, the answer is Android. In contrast, iOS is
a proprietary mobile operating system, and Windows and macOS are proprietary
desktop operating systems.

\section{Exam retrieval and common confusions}
Five distinctions prevent most errors on this topic.
\begin{enumerate}
  \item The test of open source is the \textbf{source code}, not the price. An
    open-source licence lets users \textbf{view, change, and share} the source
    (PYQ-26).
  \item \textbf{Freeware} is free of \textbf{cost}; \textbf{free software} is
    about \textbf{freedom}. Do not treat the two as the same (TRAP-23).
  \item \textbf{Open source}, \textbf{freeware}, and \textbf{shareware} are
    three different things: one shares the source, one is free of cost, one is
    a paid trial (TRAP-23).
  \item The open-source \textbf{mobile operating system} is \textbf{Android}
    (PYQ-24); iOS is proprietary.
  \item \textbf{GPL} is copyleft (derivatives stay open); \textbf{MIT},
    \textbf{Apache 2.0}, and \textbf{BSD} are permissive.
\end{enumerate}

\begin{example}
An item asks, ``Which of the following is an open-source mobile operating
system?'' Trace the options: Android is open source (built on the Linux
kernel), while iOS, Windows, and macOS are proprietary. The answer is Android.
If the stem instead asks what an open-source licence permits, the answer is the
right to view, change, and share the source code.
\end{example}

\subsection*{Revision at a glance}
\begin{center}
\begin{tabular}{p{0.20\textwidth}p{0.32\textwidth}p{0.38\textwidth}}
\toprule
\textbf{Term} & \textbf{What it is} & \textbf{What it is not} \\
\midrule
Open source & Source code is public and licensed for view/change/share & Not merely free of cost \\
Proprietary & Source kept secret; only the executable is distributed & Not necessarily paid \\
Freeware & Free of cost & Not necessarily open source \\
Free software & Freedom to run, study, change, share & Not about zero price \\
Public domain & No copyright owner & Not the same as open source \\
Shareware & Free trial, then pay & Not free forever \\
Commercial & Sold for profit & Not necessarily proprietary or open \\
Copyleft (GPL) & Derivatives must stay open & Not permissive \\
Permissive (MIT) & Reuse with few conditions & Not copyleft \\
\bottomrule
\end{tabular}
\end{center}

\subsection*{Exercises}
\begin{enumerate}[label=\arabic*.]
  \item Define source code and explain how it differs from object code.
  \item State the three things an open-source licence typically permits a user
    to do with the source code.
  \item Distinguish freeware from free software.
  \item Explain what public domain software is and how it differs from
    open-source software.
  \item Name one copyleft licence and one permissive licence, and state the
    main condition of each.
\end{enumerate}

\subsection*{Solutions}
\begingroup\small
\begin{enumerate}[label=\arabic*.]
  \item Source code is the human-readable set of instructions written by a
    programmer in a programming language. It is converted by a compiler or
    interpreter into object code, the machine-readable form the computer
    actually runs.
  \item View the source code, modify it, and share (redistribute) it, subject
    to the licence terms.
  \item Freeware is free of cost but its source code is usually not provided.
    Free software is about freedom: the user may study, change, and share the
    software; ``free'' means liberty, not price.
  \item Public domain software has no copyright owner and carries no
    restrictions at all. Open-source software is still copyrighted, but its
    licence grants the freedom to view, change, and share the source.
  \item Copyleft: the GNU GPL, which requires modified versions to be released
    under the same licence. Permissive: the MIT License (or Apache 2.0 / BSD),
    which allows reuse with few conditions, such as keeping the licence
    notice.
\end{enumerate}
\endgroup
\end{document}
